\section{Sources des illustrations}\label{annexe-illustrations}

Les fac-similés reproduisent les photographies du carnet prises le 26~novembre 2025, redressées et recadrées au bord du feuillet, l'éclairage et le ton du papier égalisés d'une page à l'autre (l'écriture n'est pas retouchée), et les numérisations des pièces annexes, redressées et recadrées au bord du papier ; les cartes ont été dessinées pour la présente édition (sources indiquées sous chacune d'elles). Les photographies, lettres de guerre, notes et feuillets autographes (\textit{Feuillets de Route}) appartiennent aux archives familiales des descendants d'Hippolyte Combes, communiquées par Pierre Combes et Pierre-Alain Hauseux. Les images d'époque, toutes dans le domaine public ou sous licence libre, sont reproduites avec les mentions demandées par les institutions qui les conservent.

{\footnotesize
\setlength{\tabcolsep}{3pt}
\renewcommand{\arraystretch}{1.12}
\begin{longtable}{@{}>{\raggedright\arraybackslash}p{50mm}>{\raggedright\arraybackslash}p{88mm}@{}}
\toprule
\textbf{Image} & \textbf{Auteur, source et licence}\\
\midrule\endhead
\bottomrule\endfoot
Hippolyte Combes en uniforme, coiffé du képi, pendant la guerre (portrait de studio) & Archives familiales (photographies militaires).\\
\textit{Bailleul (Nord), Grand-Place} — vers 1900--1909 (carte ayant circulé en 1909) & Photographe anonyme ; éditeur au trèfle à quatre feuilles (Maurice Tesson, Limoges). Collection Jpcuvelier. Source : \url{https://commons.wikimedia.org/wiki/File:Nord_-_28_-_Bailleul_%E2%80%94_H%C3%B4tel_de_Ville.jpg} ; licence : domaine public (PD-old-70-expired).\\
\textit{Bailleul (Nord), Grand-Place} — 15~mai 1917 & Henry Armytage Sanders, photographe officiel néo-zélandais ; Alexander Turnbull Library, National Library of New Zealand, réf.~1/2-013396-G. Source : \url{https://commons.wikimedia.org/wiki/File:New_Zealand_troops_march_through_Bailleul,_France_(21610313216).jpg} ; licence : «~No known copyright restrictions~» (Flickr-no known copyright restrictions, National Library of New Zealand, Flickr Commons) ; photographie officielle néo-zélandaise de 1917, \textit{Crown copyright} (100~ans) expiré. Crédit : Alexander Turnbull Library, Wellington, 1/2-013396-G.\\
Hippolyte Combes en uniforme, tête nue, pendant la guerre (portrait de studio) & Archives familiales (photographies militaires).\\
\textit{Saint-Pol-sur-Ternoise (Pas-de-Calais)} — 9~mars 1916 & Emmanuel-Louis Mas (Section photographique de l'armée) ; La contemporaine (Nanterre), fonds des «~Albums Valois~». Source : \url{https://commons.wikimedia.org/wiki/File:D%C3%A9fil%C3%A9_du_108e_RIT_dans_une_rue_de_Saint-Pol,_9_mars_1916.jpg} (notice d'origine : \url{https://argonnaute.parisnanterre.fr/ark:/14707/a011432125257Rw8Pmx}) ; licence : Licence Ouverte / Open Licence (Etalab), déclarée par La contemporaine pour les Albums Valois — réutilisation libre, y compris commerciale, avec mention de la source. Crédit : «~La contemporaine, Nanterre — Albums Valois (opérateur E. Mas)~».\\
\textit{Lucheux (Somme)} — années 1890--1900 (datation BnF : années 1890) & Photographe non identifié ; recueil attribué par la BnF à «~A. Karl~» (donateur) ; BnF, département des Estampes et de la photographie, btv1b10500747g (f.~340). Source : \url{https://commons.wikimedia.org/wiki/File:France-Album_-_fascicules_55-70_-_dessins_originaux_A._Karl_-_btv1b10500747g_(340_of_605).jpg} ; licence : domaine public (PD-France, PD-US-expired). Mention demandée : «~Source gallica.bnf.fr / BnF~».\\
Marguerite Calmels à vingt-cinq ans (1917) & Archives familiales (portraits et lieux de famille).\\
\textit{Paris, gare du Nord} — juillet 1915 (diffusion Meurisse : 9~juillet 1915) & Agence de presse Meurisse (photographe non identifié) ; BnF, département des Estampes et de la photographie, btv1b90437920 (Meurisse, 57527). Source : \url{https://commons.wikimedia.org/wiki/File:Arriv%C3%A9e_des_permissionnaires_%C3%A0_la_gare_du_Nord_-_btv1b90437920.jpg} ; licence : domaine public (PD-France, PD-1996, PD-US-expired : photographie de presse anonyme publiée en 1915). Mention demandée : «~Source gallica.bnf.fr / BnF~».\\
\textit{Amiens (Somme)} — dessin exécuté en 1916 (Bone arrive au front le 16~août 1916) ; lithographie publiée en 1917 & Sir (David) Muirhead Bone (1876--1953), premier artiste de guerre officiel britannique ; Royal Albert Memorial Museum \& Art Gallery (Exeter), inv.~50/1918/1/9. Source : \url{https://commons.wikimedia.org/wiki/File:Amiens_Cathedral_-_Sir_Muirhead_Bone_-_50_1918_1_9.jpg} ; licence : domaine public (PD-old-70 : auteur mort en 1953) ; numérisation versée sous CC0 par le musée (CC0). Crédit d'usage : Royal Albert Memorial Museum \& Art Gallery, Exeter.\\
\textit{Albert (Somme)} — vers 1916 (Archives of Ontario : «~ca. 1916~») & Photographe officiel canadien non identifié ; Archives of Ontario, C~224-8, n\textsuperscript{o}~I0004775. Source : \url{https://commons.wikimedia.org/wiki/File:Falling_statue_on_the_Albert_Cathedral_Tower_(I0004775).jpg} ; licence : domaine public — PD-US-expired + PD-Canada. Crédit : Archives of Ontario.\\
\textit{Albert (Somme), place devant la basilique} — vers 1916--1917 (Archives of Ontario : «~ca. 1917~» ; la statue est encore suspendue, donc avant avril 1918) & Photographe officiel canadien non identifié (Canadian War Records Office) ; album «~Canadian Expeditionary Force photograph albums~», Archives of Ontario, C~224-8, n\textsuperscript{o}~I0004828. Source : \url{https://commons.wikimedia.org/wiki/File:Albert_Cathedral_damaged_by_the_war_(I0004828).jpg} ; licence : domaine public — PD-US-expired + PD-Canada (photographie canadienne antérieure à 1949 ; publiée avant 1931). Crédit : Archives of Ontario.\\
\textit{Environs de Contay (Somme)} — septembre 1916 & Photographe inconnu (photographie de presse) ; Nationaal Archief (La Haye), fonds Spaarnestad, inv.~22057, n\textsuperscript{o}~22057\_041. Source : \url{https://commons.wikimedia.org/wiki/File:A_first_aid_lorry_with_a_convoy_near_Contay,_September_1916,_Bestanddeelnr_22057_041.tif} ; licence : CC0 1.0 (dédicace au domaine public par le Nationaal Archief ; «~Auteursrechthebbende : Publiek Domein~»).\\
\textit{Route de Saint-Pol-sur-Ternoise à Arras (Pas-de-Calais)} — février 1917 & John Warwick Brooke ; Imperial War Museums, Q~4683. Source : \url{https://commons.wikimedia.org/wiki/File:The_British_Army_on_the_Western_Front,_1914-1918_Q4683.jpg} ; licence : domaine public (PD-UKGov). Crédit : Imperial War Museums (Q~4683).\\
\textit{Bertrancourt (Somme), village voisin d'Acheux-en-Amiénois (environ 3~km)} — 16~avril 1918 & Henry Armytage Sanders, photographe officiel néo-zélandais ; Alexander Turnbull Library, National Library of New Zealand. Source : \url{https://commons.wikimedia.org/wiki/File:A_New_Zealand_Divisional_YMCA_building_on_the_Somme_(21474583676).jpg} ; licence : «~No known copyright restrictions~» (Flickr-no known copyright restrictions, National Library of New Zealand) ; photographie officielle néo-zélandaise, \textit{Crown copyright} expiré. Crédit : Alexander Turnbull Library, Wellington.\\
\textit{Millau (Aveyron), ganterie Jonquet} — années 1910 & Photographe et éditeur non identifiés (carte postale ; reproduction d'après Delcampe). Source : \url{https://commons.wikimedia.org/wiki/File:Millau_Gant_Jonquet_Salle_des_fourrures.jpg} ; licence : domaine public (PD-old-1923 sur Commons : œuvre anonyme publiée avant 1923).\\
\textit{Toulouse, place Saint-Sernin} — 28~avril 1916 & Auguste Léon, opérateur des Archives de la Planète (mission «~1916 – Provinces françaises~», Jean Brunhes, Auguste Léon et Georges Chevalier) ; Département des Hauts-de-Seine, musée Albert-Kahn, A~8326. Source : \url{https://commons.wikimedia.org/wiki/File:La_basilique_Saint-Sernin,_vue_depuis_la_place_Saint-Sernin,_rue_du_Taur_-_A8326.jpg} (notice du musée : \url{https://collections.albert-kahn.hauts-de-seine.fr/document/la-basilique-saint-sernin-vue-depuis-la-place-saint-sernin-rue-du-taur/617a7a43cf8b8968b3381b6f}) ; licence : CC0 1.0 (dédicace au domaine public par le musée départemental Albert-Kahn ; l'opérateur est mort en 1942). Crédit d'usage : «~Musée départemental Albert-Kahn, Département des Hauts-de-Seine~».\\
\textit{Millau (Aveyron), le Tarn et le moulin du Pont-Vieux} — entre 1890 et 1914 & Émile Duchemin (1862--1914) ; BnF, département des Estampes et de la photographie, btv1b106645153. Source : \url{https://commons.wikimedia.org/wiki/File:Aveyron_-_btv1b106645153.jpg} ; licence : domaine public (PD-France, PD-old-auto-expired [auteur mort en 1914], PD-US-expired). Mention demandée : «~Source gallica.bnf.fr / BnF~».\\
Mariage d'Hippolyte Combes et de Marguerite Calmels, Millau, 25~octobre 1917 & Archives familiales (état civil et actes familiaux).\\
\textit{Puchevillers (Somme)} — septembre 1916 & Aviation allemande, observateur «~Lt. d. R. Münz~» ; Landesarchiv Baden-Württemberg, Generallandesarchiv Karlsruhe, 456 F 7 Nr.~114 Foto 14. Source : \url{https://commons.wikimedia.org/wiki/File:Foto_des_Bahnhofs_von_Puchevillers_mit_Kriegsrampe_und_Lazarettanlage_(Luftbild)_-_LABW_-_Generallandesarchiv_Karlsruhe_456_F_7_Nr._114_Foto_14.jpg} (permalien de l'archive : \url{https://www.landesarchiv-bw.de/plink/?f=4-1701618-1}) ; licence : CC BY 4.0. Crédit obligatoire : «~Landesarchiv Baden-Württemberg, Generallandesarchiv Karlsruhe 456 F 7 Nr.~114 Foto 14 (CC BY 4.0)~».\\
\textit{Près de Pys (Somme), ancien champ de bataille de l'Ancre} — mars 1917 & John Warwick Brooke ; Imperial War Museums, Q~4997. Source : \url{https://commons.wikimedia.org/wiki/File:The_German_Withdrawal_To_the_Hindenburg_Line,_March-april_1917_Q4997.jpg} ; licence : domaine public (PD-UKGov). Crédit : Imperial War Museums (Q~4997).\\
Millau, la maison Combes (photographie récente) & Archives familiales (portraits et lieux de famille).\\
Hippolyte Combes à vingt et un ans (1897) & Archives familiales (portraits et lieux de famille).\\
Hippolyte Combes et ses camarades canonniers à l'école d'artillerie de Toulon (avril 1915) & Archives familiales (photographies militaires).\\
Menu du repas de noces, 25~octobre 1917 & Archives familiales (état civil et actes familiaux).\\
Les Andelys, l'École militaire préparatoire, chambre d'Hippolyte Combes (1918–1919) & Archives familiales (photographies militaires).\\
Hippolyte et Marguerite Combes et quatre de leurs enfants (milieu des années 1920) & Archives familiales (photographies de famille).\\
Hippolyte Combes à Cambridge (années 1930) & Archives familiales (portraits et lieux de famille).\\
\textit{Albert (Somme), dépôt de ravitaillement au bord de la route} — mars 1917 & John Warwick Brooke, photographe officiel britannique ; Imperial War Museums, Q~4828. Source : \url{https://commons.wikimedia.org/wiki/File:The_Army_Service_Corps_on_the_Western_Front,_1914-1918_Q4828.jpg} ; licence : domaine public (PD-UKGov : photographie de la Couronne britannique antérieure à 1957, droits expirés). Crédit : Imperial War Museums (Q~4828).\\
\textit{Route Amiens – Saint-Quentin, de Foucaucourt vers Brie (Somme)} — mars 1917 & Ernest Brooks, photographe officiel britannique ; Imperial War Museums, Q~1863. Source : \url{https://commons.wikimedia.org/wiki/File:Sherwood_Foresters_following_up_Germans_near_Brie_March_1917_IWM_Q_1863.jpg} ; licence : domaine public (PD-UKGov). Crédit : Imperial War Museums (Q~1863).\\
\end{longtable}}
